% TOP_PROBLEM: {"id": 30007337, "title": "Full-boundary recovery of both smooth isotropic elastic moduli", "category_id": 9, "category": {"id": 9, "name": "pde", "display_name": "Partial Differential Equations", "description": "PDEs and their applications in physics and geometry.", "slug": "pde", "order_index": 9, "created_at": "2026-07-31T15:26:25.671Z"}, "status": "open"}
% FIRST_POSTED: 2026-10-06
% ATTEMPT_STATUS: solved
% OPENAI_CLAIM: {"schema_version":1,"source_repo":"https://github.com/openai/math","source_commit":"adc7f1241b42e322a6451854ab7e4b4c146bf78a","release_date":"2026-10-06","match":"full","resolution":"proved","independently_reviewed":false,"second_pass":"agreed","families":[{"family":"372","title":"Global uniqueness in smooth isotropic elasticity","subject":"Partial differential equations","manuscripts":[{"title":"Global Uniqueness for the Smooth Isotropic Elasticity Inverse Problem","pdf_path":"preprints/Global-Uniqueness-for-the-Smooth-Isotropic-Elasticity-Inverse-Problem-September-24-2026/article.pdf","readme_path":"preprints/Global-Uniqueness-for-the-Smooth-Isotropic-Elasticity-Inverse-Problem-September-24-2026/README.md","date":"2026-09-24","theorem_ref":"preprints/Global-Uniqueness-for-the-Smooth-Isotropic-Elasticity-Inverse-Problem-September-24-2026/build/sections/introduction.tex:31 (thm:main)"}],"lean":{"doc_path":"lean/docs/372.md","comparators":["lean/ComparatorChallenges/ElasticityUniqueness.json","lean/ComparatorChallenges/ElasticityUniqueness.lean"],"declarations":[],"covers_main_theorem":true},"reasoning_trace":null}]}
\documentclass[11pt]{article}
\usepackage[margin=25mm]{geometry}
\usepackage{fontspec}
\setmainfont{Latin Modern Roman}
\usepackage{amsmath,amssymb,mathtools}
\usepackage{unicode-math}
\setmathfont{Latin Modern Math}
\usepackage{enumitem}
\usepackage{xurl,hyperref,bookmark}
\hypersetup{colorlinks=true,urlcolor=blue}
\setcounter{secnumdepth}{0}
\setlength{\parindent}{0pt}
\setlength{\parskip}{5pt}
\setlength{\emergencystretch}{3em}
\begin{document}
\section*{Full-boundary recovery of both smooth isotropic elastic moduli}

\subsection{OpenAI's claim}
\textit{Smooth isotropic elasticity uniqueness}: The announced uniqueness statement recovers both smooth Lamé coefficients from full boundary displacement and traction data on a bounded smooth three-dimensional domain when $\mu>0$ and $3\lambda+2\mu>0$. It is a uniqueness assertion under these positivity hypotheses.

\subsection{Outline of the argument}
\textit{Smooth isotropic elasticity uniqueness}: Boundary determination supplies matching jets and common exterior coefficient extensions. Equality of displacement and traction permits physical solutions in the two media to be glued across the boundary, including their actual divergences. This exact exterior comparison is propagated through the interior system, identifying both moduli rather than only an auxiliary transformed coefficient.

\subsection{Scope relative to this problem}
\textit{Smooth isotropic elasticity uniqueness}: Both Lamé coefficients are smooth real functions, with the strict isotropic positivity conditions on the closure. Full traction-map equality implies equality of both functions throughout the domain; no reconstruction or partial-boundary conclusion is added.

With the hypotheses and conventions above, the claimed result addresses the whole mathematical assertion on this page.

\subsection{Formal verification}
Family 372 describes Lean coverage for uniqueness of the two smooth positive-isotropic Lamé moduli from the full static traction map.

\textit{Smooth isotropic elasticity uniqueness}: The manuscript links \texttt{ElasticityUniqueness.lean}.

The documented selection includes this manuscript’s main assertion; that assertion retains the restrictions stated above.

This site has not run the Lean development. These links describe OpenAI’s selected formal statements and should be read with their stated scope.

\subsection{Status caveat}
OpenAI reports that an internal model produced the manuscripts. Their arguments have not been independently verified as proofs by this site. Any portion outside the documented formal scope may contain mistakes, and the catalogue label reports the announced mathematical claim.

\subsection{Sources}
\begin{itemize}
\item OpenAI, \emph{Global Uniqueness for the Smooth Isotropic Elasticity Inverse Problem}, 2026-09-24; citation key \texttt{OAI:Global-Uniqueness-for-the-Smooth-Isotropic-Elasticity-Inverse-Problem-September-24-2026}. \href{https://github.com/openai/math/blob/adc7f1241b42e322a6451854ab7e4b4c146bf78a/preprints/Global-Uniqueness-for-the-Smooth-Isotropic-Elasticity-Inverse-Problem-September-24-2026/article.pdf}{PDF (GitHub)}; \href{https://github.com/openai/math/raw/adc7f1241b42e322a6451854ab7e4b4c146bf78a/preprints/Global-Uniqueness-for-the-Smooth-Isotropic-Elasticity-Inverse-Problem-September-24-2026/article.pdf}{PDF (direct)}; \href{https://github.com/openai/math/blob/adc7f1241b42e322a6451854ab7e4b4c146bf78a/preprints/Global-Uniqueness-for-the-Smooth-Isotropic-Elasticity-Inverse-Problem-September-24-2026/README.md}{Manuscript page and citation}; \href{https://github.com/openai/math/blob/main/preprints/Global-Uniqueness-for-the-Smooth-Isotropic-Elasticity-Inverse-Problem-September-24-2026/article.pdf}{Latest version}.
\item Family 372: \href{https://github.com/openai/math/blob/adc7f1241b42e322a6451854ab7e4b4c146bf78a/lean/docs/372.md}{Lean scope}.
\item \href{https://github.com/openai/math/blob/adc7f1241b42e322a6451854ab7e4b4c146bf78a/lean/ComparatorChallenges/ElasticityUniqueness.json}{Comparator ElasticityUniqueness.json}.
\item \href{https://github.com/openai/math/blob/adc7f1241b42e322a6451854ab7e4b4c146bf78a/lean/ComparatorChallenges/ElasticityUniqueness.lean}{Comparator ElasticityUniqueness.lean}.
\item \href{https://github.com/openai/math/blob/adc7f1241b42e322a6451854ab7e4b4c146bf78a/CONTENTS.md}{OpenAI catalogue entry}.
\end{itemize}
\end{document}
